In \ref{chap:finetuning} a model was fine-tuned to answer a \termref{value-survey} in a similar manner to a human. That shows the model can reproduce survey answers, but not that it behaves differently when it has to make a decision. We therefore built behavioural simulations: an LLM placed in an \termref{agentic-harness}, given a work profile and a situation, and left to work the case with tools. The aim was to test whether the fine-tuned models made appreciably different decisions from the base model (\rqref{align}), and to produce trajectory summaries that the public could compare (\rqref{prefer}).

The approach was shelved. The harness was built and run, but the simulations did not produce clean, interpretable value differences between models: the models did not diverge reliably on the value-laden parts of the task, and the extra freedom of an agentic trajectory made the decisions hard to compare. The project therefore pivoted to the single-response dilemmas in \ref{chap:dilemmas}. This appendix records what was built and what is worth reusing in a follow-up.


\section{What was tried}

\subsection{Agentic harness}

Each simulation is an LLM in a ReAct-style harness \cite{yaoReActSynergizingReasoning2023}. It receives a work profile (system prompt and tool set) and a situation (a work item), and can reason, call tools, and continue until a situation-specific terminal event such as closing a case or submitting an application. Interactive situations add an LLM interlocutor with a strict persona and instructions, so the agent can ask follow-up questions. To keep the setup realistic, prompts and tools were based on public guidelines and documentation rather than hand-written value cues, and profile data was modelled on real organisational records.

\subsection{Profiles and situations}

Each profile is a plausible deployment context with its own tools; each situation is one work item within that context. We built five profiles with five situations each (\tabref{behavioural-sim-profiles}). The profiles cover high-stakes work that normally requires human judgement in healthcare, finance and government services, because these are the areas the public is most concerned about \cite{Aitrustreport2025onenz,GlobalAttitudesAI2026}. Situations had to be value-first (the key decisions have no obviously correct answer), relevant to the New Zealand public, and understandable without domain expertise.

\ctable{behavioural-sim-profiles}{Simulation profiles and situations. Bolded entries are interactive situations.}

Situations were chosen to expose value tensions between New Zealand and the countries where most frontier models are developed (the USA and China). Using the Inglehart--Welzel cultural map \cite{WVSMAP}, we identified tensions on the traditional--secular and survival--self-expression dimensions and mapped them to Schwartz values \cite{schwartzRefiningTheoryBasic2012}. We did not choose the questions where \modelref{qwen3.8-27b} diverged most from New Zealand: those questions were not significantly more divergent than the rest, so selecting them would mostly fit noise.


\section{Reviewing trajectories}

A run produces a trajectory: the setup, the model's reasoning, its tool calls and, in interactive situations, the interlocutor's replies. Trajectories were reviewed in two ways (\figref{behavioural-sim-pipeline}). A structured review used a situation-specific grading rubric with typed items (Likert, multichoice, boolean or number), completed by \judgemodel{}. This yields comparable variables for cross-model analysis and for clustering trajectory types. A holistic review asked the target model, still in character, to write a short self-review; \judgemodel{} then audited that review for fairness. The audited summary was intended for the public pairwise comparison in \ref{chap:survey}.

\cdiagram{behavioural-sim-pipeline}{The behavioural simulation pipeline. The trajectory splits into two review paths: a structured review by an external model for the cross-model analysis (RQ1), and a self-review by the target model, audited by an external judge, that feeds the public survey in the next chapter (RQ2).}


\section{Lessons for a follow-up study}

The main lesson is that agentic freedom is a double-edged sword. Realistic tools and multi-step cases make behaviour more deployment-like, but they also multiply the number of decisions on which models can differ, and most of those differences are procedural rather than value-laden. A follow-up should keep the realistic harness, environment data and terminal events, but reduce each situation to a small number of explicit, value-laden decision points; pre-test the situations on several models and keep only those where the models actually diverge on those decisions; and retain the structured rubric plus the audited self-review so that the same runs can support both cross-model analysis and public comparison.

One caveat is realism. LLMs can detect that they are in a simulation and behave differently when they believe they are being evaluated \cite{berglundTakenOutContext2023,needhamLargeLanguageModels2025,greenblattAlignmentFakingLarge2024}. We tried to reduce this by using realistic organisation data, avoiding explicit value prompts, and presenting cases as ordinary work items. A follow-up should still check the trajectories for evaluation awareness and test whether behaviour changes when the framing changes.
